\expandafter\gdef\csname SNside1\endcsname{\node[anchor=north west,inner sep=0pt,text width=48mm,font=\SideFont,align=left] (SNsidebox) at ([xshift=138mm,yshift=-402.91pt]current page.north west) {\textcolor{SNink}{\textbf{不变量 · 路径形成与查表分开}}\par\vspace{1.5mm}动态路由更新可以较慢，当前分组的查表转发仍须高速执行。题目给出``路由器收到一份分组''并不意味着它立刻运行一次全网最短路；应先使用已经形成的转发表判断下一跳与出口。};\pgfextracty{\SNsideheight}{\pgfpointdiff{\pgfpointanchor{SNsidebox}{south}}{\pgfpointanchor{SNsidebox}{north}}}\typeout{SNHEIGHT:1:0:\the\SNsideheight}}
\expandafter\gdef\csname SNside2\endcsname{\node[anchor=north west,inner sep=0pt,text width=48mm,font=\SideFont,align=left] (SNsidebox) at ([xshift=138mm,yshift=-95.00pt]current page.north west) {\textcolor{SNink}{\textbf{课程连接 · 控制像策略数据像执行}}\par\vspace{1.5mm}操作系统调度器决定策略，CPU按选定上下文执行；SDN也把规则计算与逐包执行分开。这个类比只帮助理解分工，不能据此认定控制器必须单机、必须在每份分组的实际转发路径上。};\pgfextracty{\SNsideheight}{\pgfpointdiff{\pgfpointanchor{SNsidebox}{south}}{\pgfpointanchor{SNsidebox}{north}}}\typeout{SNHEIGHT:2:1:\the\SNsideheight}}
\expandafter\gdef\csname SNside3\endcsname{\node[anchor=north west,inner sep=0pt,text width=48mm,font=\SideFont,align=left] (SNsidebox) at ([xshift=138mm,yshift=-95.00pt]current page.north west) {\textcolor{SNink}{\textbf{不变量 · 暂定距离尚可变小}}\par\vspace{1.5mm}Dijkstra的未确定点距离只是当前最好候选，选定最小点后再松弛其他点。A到C初值4，在经B的路径代价1+2出现后应降为3；算法正确性还依赖非负边权，不能省略条件。};\pgfextracty{\SNsideheight}{\pgfpointdiff{\pgfpointanchor{SNsidebox}{south}}{\pgfpointanchor{SNsidebox}{north}}}\typeout{SNHEIGHT:3:4:\the\SNsideheight}}
\expandafter\gdef\csname SNside4\endcsname{\node[anchor=north west,inner sep=0pt,text width=48mm,font=\SideFont,align=left] (SNsidebox) at ([xshift=138mm,yshift=-95.00pt]current page.north west) {\textcolor{SNink}{\textbf{最小反例 · MF零不等于未分片}}\par\vspace{1.5mm}未分片数据报可有偏移0、MF=0；一个最后分片也可有偏移97、MF=0。两个状态共享MF数值，已足以排除``看到MF为0就断定没有分片''，必须同时读取片偏移。\par\smallskip\begin{tikzpicture}[x=1mm,y=1mm,>=Stealth,node distance=5mm]
\node[draw=SNink!50,fill=SNblue,rounded corners=1mm,text width=34mm,align=center,inner sep=1.2mm,font=\fontsize{7.5}{9}\selectfont] (m0) at (0,0) {未分片};
\node[draw=SNink!50,fill=SNblue,rounded corners=1mm,text width=34mm,align=center,inner sep=1.2mm,font=\fontsize{7.5}{9}\selectfont] (m1) at (0,-11) {偏0 MF0};
\node[draw=SNink!50,fill=SNblue,rounded corners=1mm,text width=34mm,align=center,inner sep=1.2mm,font=\fontsize{7.5}{9}\selectfont] (m2) at (0,-22) {末片};
\node[draw=SNink!50,fill=SNblue,rounded corners=1mm,text width=34mm,align=center,inner sep=1.2mm,font=\fontsize{7.5}{9}\selectfont] (m3) at (0,-33) {偏97 MF0};
\draw[->,SNink] (m0) -- (m1);
\draw[->,SNink] (m2) -- (m3);
\end{tikzpicture}
\par{\fontsize{7.5}{9}\selectfont MF同为0，偏移揭示不同状态}};\pgfextracty{\SNsideheight}{\pgfpointdiff{\pgfpointanchor{SNsidebox}{south}}{\pgfpointanchor{SNsidebox}{north}}}\typeout{SNHEIGHT:4:6:\the\SNsideheight}}
\expandafter\gdef\csname SNside5\endcsname{\node[anchor=north west,inner sep=0pt,text width=48mm,font=\SideFont,align=left] (SNsidebox) at ([xshift=138mm,yshift=-95.00pt]current page.north west) {\textcolor{SNink}{\textbf{不变量 · 偏移坐标从原数据起}}\par\vspace{1.5mm}分片的偏移以原始IP数据起点为零，单位8 B，不把任何分片首部累计进去。再次分片也沿用同一原坐标，因此原片的最后子片若后面仍有原数据，MF仍应保持1。};\pgfextracty{\SNsideheight}{\pgfpointdiff{\pgfpointanchor{SNsidebox}{south}}{\pgfpointanchor{SNsidebox}{north}}}\typeout{SNHEIGHT:5:5:\the\SNsideheight}}
\expandafter\gdef\csname SNside6\endcsname{\node[anchor=north west,inner sep=0pt,text width=48mm,font=\SideFont,align=left] (SNsidebox) at ([xshift=138mm,yshift=-95.00pt]current page.north west) {\textcolor{SNink}{\textbf{错项排除 · 普通主机式有用途前提}}\par\vspace{1.5mm}普通多址子网常从地址总数减去网络和广播两个地址；主机路由/32、点到点/31等用途有其他规则。先判断题设用途，再套主机数公式，不能把二的幂减二写成所有前缀通用定律。};\pgfextracty{\SNsideheight}{\pgfpointdiff{\pgfpointanchor{SNsidebox}{south}}{\pgfpointanchor{SNsidebox}{north}}}\typeout{SNHEIGHT:6:7:\the\SNsideheight}}
\expandafter\gdef\csname SNside7\endcsname{\node[anchor=north west,inner sep=0pt,text width=48mm,font=\SideFont,align=left] (SNsidebox) at ([xshift=138mm,yshift=-95.00pt]current page.north west) {\textcolor{SNink}{\textbf{最小反例 · 更具体的出口会胜出}}\par\vspace{1.5mm}目的10.1.2.7同时落在10/8与10.1.2/24中，即使/8记录的距离更小，转发仍选/24。先比较匹配前缀长度，才处理同一前缀候选的度量；表项先后顺序也不决定这个结果。};\pgfextracty{\SNsideheight}{\pgfpointdiff{\pgfpointanchor{SNsidebox}{south}}{\pgfpointanchor{SNsidebox}{north}}}\typeout{SNHEIGHT:7:8:\the\SNsideheight}}
\expandafter\gdef\csname SNside8\endcsname{\node[anchor=north west,inner sep=0pt,text width=48mm,font=\SideFont,align=left] (SNsidebox) at ([xshift=138mm,yshift=-95.00pt]current page.north west) {\textcolor{SNink}{\textbf{不变量 · ARP先确定当前邻居}}\par\vspace{1.5mm}源主机的第一帧只需交给本地下一跳，ARP广播也只到当前广播域。即使IP最终目的在远端，当前帧目的MAC仍可为网关；普通二层交换机转发这帧时，不把源MAC换成自己。};\pgfextracty{\SNsideheight}{\pgfpointdiff{\pgfpointanchor{SNsidebox}{south}}{\pgfpointanchor{SNsidebox}{north}}}\typeout{SNHEIGHT:8:10:\the\SNsideheight}}
\expandafter\gdef\csname SNside9\endcsname{\node[anchor=north west,inner sep=0pt,text width=48mm,font=\SideFont,align=left] (SNsidebox) at ([xshift=138mm,yshift=-95.00pt]current page.north west) {\textcolor{SNink}{\textbf{错项排除 · 地址提议尚非地址生效}}\par\vspace{1.5mm}Offer提供候选地址，初始选择阶段Request表达接受哪份提议，ACK才确认租约。图中出现候选IP并不意味着它已可作客户端IP源；题目若问目的、源，还应按设问顺序填写。};\pgfextracty{\SNsideheight}{\pgfpointdiff{\pgfpointanchor{SNsidebox}{south}}{\pgfpointanchor{SNsidebox}{north}}}\typeout{SNHEIGHT:9:11:\the\SNsideheight}}
\expandafter\gdef\csname SNside10\endcsname{\node[anchor=north west,inner sep=0pt,text width=48mm,font=\SideFont,align=left] (SNsidebox) at ([xshift=138mm,yshift=-95.00pt]current page.north west) {\textcolor{SNink}{\textbf{错项排除 · 差错报告不是可靠重传}}\par\vspace{1.5mm}ICMP把某次网络问题告知原源端，不负责替所有应用重发丢失分组。差错之间还有限制以防反馈放大；历史源点抑制已被废弃，不能把旧类型识别题的答案当成当前拥塞处理建议。};\pgfextracty{\SNsideheight}{\pgfpointdiff{\pgfpointanchor{SNsidebox}{south}}{\pgfpointanchor{SNsidebox}{north}}}\typeout{SNHEIGHT:10:12:\the\SNsideheight}}
\expandafter\gdef\csname SNside11\endcsname{\node[anchor=north west,inner sep=0pt,text width=48mm,font=\SideFont,align=left] (SNsidebox) at ([xshift=138mm,yshift=-204.47pt]current page.north west) {\textcolor{SNink}{\textbf{错项排除 · 原下一跳变差也得更新}}\par\vspace{1.5mm}本机原路由经B到N为3跳，B现在通告5跳，经B候选就变为6。继续保留旧3会把已失效的信息当成更优路线；只有来自不同下一跳的更差候选，才可按通常规则忽略。};\pgfextracty{\SNsideheight}{\pgfpointdiff{\pgfpointanchor{SNsidebox}{south}}{\pgfpointanchor{SNsidebox}{north}}}\typeout{SNHEIGHT:11:14:\the\SNsideheight}}
\expandafter\gdef\csname SNside12\endcsname{\node[anchor=north west,inner sep=0pt,text width=48mm,font=\SideFont,align=left] (SNsidebox) at ([xshift=138mm,yshift=-95.00pt]current page.north west) {\textcolor{SNink}{\textbf{不变量 · BGP策略先于简单距离}}\par\vspace{1.5mm}AS\_PATH帮助抑制AS环并参与选路，真实决策还受本地优先级等属性约束。只有题目明确其余策略属性相同，才可按路径长度及到出口的域内代价进一步比较，不能固定背成``AS少就选''。};\pgfextracty{\SNsideheight}{\pgfpointdiff{\pgfpointanchor{SNsidebox}{south}}{\pgfpointanchor{SNsidebox}{north}}}\typeout{SNHEIGHT:12:15:\the\SNsideheight}}
\expandafter\gdef\csname SNside13\endcsname{\node[anchor=north west,inner sep=0pt,text width=48mm,font=\SideFont,align=left] (SNsidebox) at ([xshift=138mm,yshift=-95.00pt]current page.north west) {\textcolor{SNink}{\textbf{课程连接 · 成员管理与转发树分工}}\par\vspace{1.5mm}IGMP像维护当前接收者集合，组播路由机制再按网络结构组织复制分支。知道本地有人加入，并不等于已经知道远端所有路径；TTL只能限制错误环流寿命，不能替代合理的树结构。};\pgfextracty{\SNsideheight}{\pgfpointdiff{\pgfpointanchor{SNsidebox}{south}}{\pgfpointanchor{SNsidebox}{north}}}\typeout{SNHEIGHT:13:16:\the\SNsideheight}}
\expandafter\gdef\csname SNside14\endcsname{\node[anchor=north west,inner sep=0pt,text width=48mm,font=\SideFont,align=left] (SNsidebox) at ([xshift=138mm,yshift=-95.00pt]current page.north west) {\textcolor{SNink}{\textbf{不变量 · 内外层各有自己的目的}}\par\vspace{1.5mm}移动节点的归属地址保持身份，转交地址描述当前位置。隧道外层把分组送到当前访问网络，内层仍保留原目的；将二者画成两个封装层，就不会误以为每一跳都把原目的IP改写。};\pgfextracty{\SNsideheight}{\pgfpointdiff{\pgfpointanchor{SNsidebox}{south}}{\pgfpointanchor{SNsidebox}{north}}}\typeout{SNHEIGHT:14:17:\the\SNsideheight}}
\expandafter\gdef\csname SNside15\endcsname{\node[anchor=north west,inner sep=0pt,text width=48mm,font=\SideFont,align=left] (SNsidebox) at ([xshift=138mm,yshift=-154.13pt]current page.north west) {\textcolor{SNink}{\textbf{课程连接 · 成员管理与转发树分工}}\par\vspace{1.5mm}IGMP像维护当前接收者集合，组播路由机制再按网络结构组织复制分支。知道本地有人加入，并不等于已经知道远端所有路径；TTL只能限制错误环流寿命，不能替代合理的树结构。};\pgfextracty{\SNsideheight}{\pgfpointdiff{\pgfpointanchor{SNsidebox}{south}}{\pgfpointanchor{SNsidebox}{north}}}\typeout{SNHEIGHT:15:16:\the\SNsideheight}}
\expandafter\gdef\csname SNside16\endcsname{\node[anchor=north west,inner sep=0pt,text width=48mm,font=\SideFont,align=left] (SNsidebox) at ([xshift=138mm,yshift=-95.00pt]current page.north west) {\textcolor{SNink}{\textbf{不变量 · 内外层各有自己的目的}}\par\vspace{1.5mm}移动节点的归属地址保持身份，转交地址描述当前位置。隧道外层把分组送到当前访问网络，内层仍保留原目的；将二者画成两个封装层，就不会误以为每一跳都把原目的IP改写。};\pgfextracty{\SNsideheight}{\pgfpointdiff{\pgfpointanchor{SNsidebox}{south}}{\pgfpointanchor{SNsidebox}{north}}}\typeout{SNHEIGHT:16:17:\the\SNsideheight}}
\expandafter\gdef\csname SNside17\endcsname{\node[anchor=north west,inner sep=0pt,text width=48mm,font=\SideFont,align=left] (SNsidebox) at ([xshift=138mm,yshift=-243.72pt]current page.north west) {\textcolor{SNink}{\textbf{课程连接 · 成员管理与转发树分工}}\par\vspace{1.5mm}IGMP像维护当前接收者集合，组播路由机制再按网络结构组织复制分支。知道本地有人加入，并不等于已经知道远端所有路径；TTL只能限制错误环流寿命，不能替代合理的树结构。};\pgfextracty{\SNsideheight}{\pgfpointdiff{\pgfpointanchor{SNsidebox}{south}}{\pgfpointanchor{SNsidebox}{north}}}\typeout{SNHEIGHT:17:16:\the\SNsideheight}}
\expandafter\gdef\csname SNside18\endcsname{\node[anchor=north west,inner sep=0pt,text width=48mm,font=\SideFont,align=left] (SNsidebox) at ([xshift=138mm,yshift=-95.00pt]current page.north west) {\textcolor{SNink}{\textbf{辨析 · 分类与CIDR}}\par\vspace{1.5mm}分类编址按地址前缀判断类别；CIDR按明确前缀长度解释，不能混用默认掩码。};\pgfextracty{\SNsideheight}{\pgfpointdiff{\pgfpointanchor{SNsidebox}{south}}{\pgfpointanchor{SNsidebox}{north}}}\typeout{SNHEIGHT:18:18:\the\SNsideheight}}
\expandafter\gdef\csname SNside19\endcsname{\node[anchor=north west,inner sep=0pt,text width=48mm,font=\SideFont,align=left] (SNsidebox) at ([xshift=138mm,yshift=-95.00pt]current page.north west) {\textcolor{SNink}{\textbf{方法 · 先需求后对齐}}\par\vspace{1.5mm}把主机与网关需求合计，再选择足够大的2次幂地址块，并检查边界对齐。};\pgfextracty{\SNsideheight}{\pgfpointdiff{\pgfpointanchor{SNsidebox}{south}}{\pgfpointanchor{SNsidebox}{north}}}\typeout{SNHEIGHT:19:19:\the\SNsideheight}}
\expandafter\gdef\csname SNside20\endcsname{\node[anchor=north west,inner sep=0pt,text width=48mm,font=\SideFont,align=left] (SNsidebox) at ([xshift=138mm,yshift=-216.43pt]current page.north west) {\textcolor{SNink}{\textbf{条件 · 转发表还需递归解析}}\par\vspace{1.5mm}BGP的NEXT\_HOP未必是本地邻居，要通过IGP确定实际出口。};\pgfextracty{\SNsideheight}{\pgfpointdiff{\pgfpointanchor{SNsidebox}{south}}{\pgfpointanchor{SNsidebox}{north}}}\typeout{SNHEIGHT:20:21:\the\SNsideheight}}
\expandafter\gdef\csname SNside21\endcsname{\node[anchor=north west,inner sep=0pt,text width=48mm,font=\SideFont,align=left] (SNsidebox) at ([xshift=138mm,yshift=-95.00pt]current page.north west) {\textcolor{SNink}{\textbf{辨析 · 自动配置有多种方式}}\par\vspace{1.5mm}IPv6既支持无状态自动配置，也支持DHCPv6，接口标识符不必来自MAC。};\pgfextracty{\SNsideheight}{\pgfpointdiff{\pgfpointanchor{SNsidebox}{south}}{\pgfpointanchor{SNsidebox}{north}}}\typeout{SNHEIGHT:21:22:\the\SNsideheight}}
\AddToShipoutPictureFG{\begin{tikzpicture}[remember picture,overlay]\ifcsname SNside\thepage\endcsname\csname SNside\thepage\endcsname\fi\end{tikzpicture}}